\chapter{Methodology}
\label{ch:methodology}

Both research tracks run through one fixed measurement apparatus. This
chapter defines it once: the data representation
(\cref{sec:panel}), the supervision label (\cref{sec:label}), the split
protocol and its leakage controls (\cref{sec:splits}), the model and
training objective (\cref{sec:training}), the four research gates
(\cref{sec:gates}), and the attention A/B that carries the core thesis
test (\cref{sec:anchors}). \Cref{sec:skill} adds the forecast-skill
apparatus used after the energy track was reframed from alpha research to
forecasting. Every experiment in \cref{ch:equity,ch:energy} reports
against this apparatus without modification; where a design choice was
made in response to an observed failure, the chapter that reports the
failure also documents the change.

\section{The canonical panel}
\label{sec:panel}

All data enters the pipeline as a panel indexed by \emph{(time, entity)}:
\emph{(date, symbol)} on the equity track, \emph{(timestamp, bidding
zone)} on the energy track. The entities are the nodes a graph later
connects. A \emph{snapshot} is the cross-section of one timestamp --- the
unit that graph propagation operates on; no operation ever sees the panel
as a whole time series.

Node features are deliberately not raw market data. Each track already
maintains a registry of island factors --- alphas that score a node from
its own history alone, ten on the equity track and eight on the energy
track. Their cross-sectionally ranked values form the feature matrix.
This choice pins the research question: any value the graph model adds
must come from relating nodes, because the per-node information is
the same information the island baselines consume. Concretely, the
equity ten are functions of each stock's own prices and volumes alone
(reversal, medium-horizon momentum, intraday range position, volume
shocks, overnight-gap behaviour); nothing exogenous enters that track.
The energy eight mix each zone's price dynamics with its published
fundamentals, such as residual load, wind and solar forecast terms,
demand surprise, and the gas spark spread. Which signals enter matters
less than the rule that they enter everywhere: island baselines,
anchors, and GAT consume this identical matrix, so a relational model
can hold no information advantage over the baselines it is measured
against. Missing values are
filled with the per-snapshot median, which in rank space is the
no-information position; a node with no data is indistinguishable from a
node ranked exactly in the middle.

\section{Supervision label}
\label{sec:label}

The label is the forward return over a fixed horizon $k$, standardised
per snapshot so that the model learns to rank the cross-section rather
than to predict absolute magnitudes. On the equity track it is
$r_{t,k} = p_{t+k}/p_t - 1$ with $k$ = 5 trading days, z-scored within
each snapshot. On the energy track it is
$r_{t,k} = (p_{t+k}-p_t)\,/\max(|p_t|, 20)$ with $k$ = 24 hours, clipped
to $\pm 0.8$ and z-scored; the denominator floor is unavoidable because
European power prices cross zero and turn negative, so a plain relative
return divides by numbers arbitrarily close to zero.

The label is supervision only. It is computed from prices at $t+k$ and
never appears as a feature; the last $k$ snapshots of each entity have no
label and are excluded from training. The energy clip, introduced for
label stability, later turned out to have a second life as an
\emph{evaluation} artifact --- \cref{ch:energy} reports how clipping the
evaluation return concealed tail losses, and why the shipped evaluation
is unclipped.

\section{Split protocol and leakage controls}
\label{sec:splits}

Leakage was treated as the primary methodological risk from the start,
and the controls are structural rather than procedural: they are asserted
in code, not remembered by the experimenter.

A single in-sample/out-of-sample boundary is placed at 70\% of the
timestamps. That one date is computed once and passed explicitly to every
consumer: graph construction builds edges only from data strictly before
it, model selection sees only data inside it, and the gate evaluation
scores only data after it. The alternative --- letting each component
derive its own boundary --- failed quietly during development: three
components each assumed 70\% and happened to agree, so changing any one
ratio would have leaked without any test failing. Centralising the date
turned that silent coupling into an explicit contract.

Within the in-sample region the layout is
\[
\texttt{train} \;\big|\; \texttt{embargo}(k) \;\big|\; \texttt{valid}
\;\big|\; \texttt{embargo}(k) \;\big|\; \texttt{OOS},
\]
with the validation block embargoed on both sides. The leading
embargo is standard practice: the last training labels reach $k$ steps
forward and must not overlap the first validation features
\cite{lopezdeprado2018afml}. The trailing embargo is easy to omit and
was nearly omitted here: validation labels also reach $k$ steps forward,
so without it, best-epoch selection would be scored partly on returns
that belong to the out-of-sample window. A dedicated assertion
(\texttt{is\_constrained\_split}) rejects any split violating this
layout.

Two automated controls argue that the pipeline as a whole produces no
false positives. A negative control shuffles the labels and asserts that
validation IC collapses to zero; a positive control plants a synthetic
signal and asserts that training recovers it. Both run in the test suite
on every change to the leakage-critical code, which is deliberately pure
pandas --- testable without a GPU or the torch dependency.

\section{Model and training objective}
\label{sec:training}

The relational model is a stacked graph attention network
\cite{velickovic2018gat} with two layers, implemented on PyTorch
Geometric \cite{fey2019pyg}. (The later forecasting work of
\cref{ch:energy} uses the GATv2 variant \cite{brody2022gatv2}; the alpha
track predates it in the codebase and keeps the original operator.) The
model maps one snapshot's node features and topology to one score per
node; scores across snapshots form the relational composite factor.

Training minimises an IC loss --- the negative Pearson correlation
between predictions and labels within each snapshot --- a differentiable
stand-in for the evaluation metric, the per-snapshot rank-IC (the
Spearman correlation of predictions and labels, reported throughout as
the OOS IC). An MSE variant is retained for comparison;
\cref{ch:equity} reports the probe that motivated the switch, in which
the IC loss recovered a planted signal far better than MSE. Model
selection is by best validation IC with epoch checkpointing: the weights
deployed are those of the best validation epoch, not the last one. This
detail matters in practice --- on the first real-data run the validation
IC peaked mid-training and then decayed to roughly half its peak, so
returning the final epoch would have discarded most of the signal.

Two deployment modes are evaluated throughout. \emph{Single} fits once on
the in-sample region and scores the entire out-of-sample window with
frozen weights. \emph{Walk-forward} refits at every 63-snapshot boundary
through the out-of-sample window, each refit using only data whose labels
predate its boundary, with the same embargoed layout per fold. Every
walk-forward score therefore comes from a model that would have been
trainable at that date in deployment. All headline runs are CPU; a GPU
benchmark found no speedup at this graph size (about 50 nodes).

\section{The four research gates}
\label{sec:gates}

A relational composite is only interesting if it survives comparison
against the island factors it was built from. Four gates, all scored on
the out-of-sample slice, formalise that comparison
(\cref{tab:gates}).

\begin{table}[htbp]
  \centering
  \caption{The four research gates. The composite is compared against
  the single island alphas (ten equity, eight energy), out-of-sample.}
  \label{tab:gates}
  \begin{tabular}{lll}
    \toprule
    Gate & Definition & Pass threshold \\
    \midrule
    Value-added & composite OOS Sharpe vs.\ best single & composite $>$ max \\
    Consistency & IS/OOS IC same sign, magnitude retention & score $\geq 0.5$ \\
    Uniqueness  & max $|$Spearman$|$ vs.\ each single & $< 0.7$ \\
    Robustness  & blend: OOS coverage, IC IR, drawdown & score $\geq 0.5$ \\
    \bottomrule
  \end{tabular}
\end{table}

The value-added gate is deliberately the strictest formulation ---
beating the maximum of $N$ competitors is a higher bar than beating their
mean, and a composite can be genuinely useful while failing it. Results
therefore also report the beat-rate against the full set of singles, and
\cref{ch:equity} discusses the one gate this project never closed.

\section{The attention A/B: two no-learning anchors}
\label{sec:anchors}

A GAT that beats the island factors has not yet demonstrated that
\emph{learned attention} was the ingredient that mattered; the
improvement could come from neighbour averaging alone. Every run
therefore carries two anchors, computed from the same inputs and, where
applicable, the same topology: the \emph{island anchor}, an equal-weight
composite of the input alphas with no propagation at all; and the
\emph{uniform anchor}, the same composite averaged uniformly over the
GAT's own topology --- relational, but with nothing learned.
\Cref{fig:iso-anchors} draws the arrangement: three factor
constructions, one input matrix, one harness.

\begin{figure}[htbp]
  \centering
  \includegraphics[width=0.88\textwidth]{isolation_anchors.pdf}
  \caption{The attention A/B. All three factor constructions read the
  same ranked island-factor matrix and are scored by the same
  evaluation harness; the two dashed comparisons are the quantities
  this report is about --- what the graph adds (island anchor versus
  uniform anchor) and what learned attention adds (uniform anchor
  versus GAT, the attention value-add).}
  \label{fig:iso-anchors}
\end{figure}

The difference between the GAT's out-of-sample Sharpe and the uniform
anchor's is the \emph{attention value-add}, the quantity the core thesis
stands on. It isolates the contribution of learned attention while
holding the information set and the graph fixed. Where the four gates ask
``is the composite a good factor?'', the A/B asks the sharper question
``did learning the edge weights matter?'' --- and the two questions turn
out to have different answers in \cref{ch:equity}.

\section{Forecast-skill apparatus (energy reframe)}
\label{sec:skill}

After the energy alpha attempt concluded negatively
(\cref{ch:energy}), the energy question was reframed from ``is there a
tradeable cross-sectional signal?'' to ``does the physical
interconnector graph improve forecast accuracy?''. The measurement
apparatus changes accordingly. The metric is a skill score against
persistence,
\[
\text{skill} = 1 - \mathrm{MSE}/\mathrm{MSE}_{\text{persistence}},
\]
scored out-of-sample at $k$ = 24 hours, alongside MAE and the
cross-sectional rank-IC. Persistence (tomorrow equals today) scores zero
by construction; a model must beat carrying yesterday forward to show any
skill at all.

Candidates form a ladder in which each rung isolates one ingredient: a
seasonal-naive rung isolates the diurnal cycle; a per-node ridge
regression on physical drivers (load, wind and solar forecasts, fuel
price) isolates what a zone's own fundamentals add; a ridge with
uniform neighbour averaging isolates what the interconnector graph adds
without learning; and the GAT rungs isolate what learned attention adds
on top. The anchors of \cref{sec:anchors} reappear here in forecasting
form --- the value claim is the margin over the no-graph rung \emph{and}
over the uniform-graph rung, not the absolute skill.

Feature timing follows the day-ahead market's information structure:
load, wind, and solar forecasts for hour $t+k$ are published
before $t$ and may legitimately be used as features for $t+k$; realised
quantities (actual load, generation mix) are known only up to $t$ and
enter as point-in-time anchors. The synthetic negative control carries
over: synthetic zones are generated independently, so the test suite
asserts that the graph rung does not beat the no-graph rung on
synthetic data. Any measured graph lift can then only come from real
cross-zonal coupling.

